\subsection{自由模上的对称张量}

命题 4. --- 设 M 是自由的，且 \((e_{i})_{i\in I}\) 是 M 的一组基。

\begin{enumerate}
\def\labelenumi{(\roman{enumi})}
\tightlist
\item
  对于 \(\nu \in \mathbf{N}^{(I)}\) ，令 \(e_{\nu} = \prod_{i\in I}\gamma_{\nu_i}(e_i)\) 。则 \((e_{\nu})_{\nu \in \mathbf{N}^{(I)}}\) 是 A-模的一组基
\end{enumerate}

\(\mathbf{TS}(\mathbf{M})\) 。特别地，代数 \(\mathbf{TS}(\mathbf{M})\) 由元素族 \(\gamma_k(x)\) （\(k \in \mathbf{N}\) 且 \(x \in \mathbf{M}\) ）生成。

\begin{enumerate}
\def\labelenumi{(\roman{enumi})}
\setcounter{enumi}{1}
\tightlist
\item
  对于每个 \(p \in \mathbf{N}\) ，\(\mathbf{T}\mathbf{S}^p(\mathbf{M})\) 是 A-模 \(\mathbf{T}^p(\mathbf{M})\) 的一个直因子。
\end{enumerate}

我们使用上述注记 2 的记号。族 \((e_{\rho(1)} \otimes \ldots \otimes e_{\rho(p)})_{\rho \in \mathcal{M}}\) 是 \(\mathbf{T}^p(\mathbf{M})\) 的一组基。因此，命题 4 由公式 (7) 及下列引理（取 \(\mathbf{H} = \mathfrak{S}_p\) 与 \(\mathbf{U} = \mathbf{T}^p(\mathbf{M})\) 应用）得出。

引理 1. --- 设 H 为有限群，U 为左 A{[}H{]}-模。假设 A-模 U 有一个基 B，它在 H 于 U 中的运算下保持稳定，并设 \(\Omega = \mathrm{B} / \mathrm{H}\) 。对于每个 \(\omega \in \Omega\) ，令 \(u_{\omega} = \sum b\) ；则

\begin{enumerate}
\def\labelenumi{(\roman{enumi})}
\item
  \((u_{\omega})_{\omega \in \Omega}\) 是 A-模 \(\mathbf{U}^{\mathrm{H}}\) 的一组基。
\item
  对于每个 \(\omega \in \Omega\) ，令 \(v_{\omega}\) 为 \(\omega\) 的一个点；令 \(\omega' = \omega - \{v_{\omega}\}\) 且 \(\mathbf{B}' = \bigcup_{\omega \in \Omega} \omega'\) ，则 \(\mathbf{B}'\) 是 \(\mathbf{U}^{\mathrm{H}}\) 在 \(\mathbf{U}\) 中的一个补空间的一组基。
\end{enumerate}

所有 \(u_{\omega}\) （\(\omega \in \Omega\) ）组成的集合与 \(B'\) 的并是 \(U\) 的一组基。若 \(U' = \sum_{\omega \in \Omega} Au_{\omega}\) 且 \(U'' = \sum_{b \in B'} Ab\) ，则我们有 \(U = U' \oplus U''\) 。另一方面，

我们有 \(u_{\omega} \in \mathbf{U}^{\mathrm{H}}\) （对所有 \(\omega \in \Omega\) ），因此 \(\mathbf{U}' \subset \mathbf{U}^{\mathrm{H}}\) 。最后，设 \((\alpha_b)_b \in B\) 是 \(A\) 的一个具有有限支撑的元素族，并令 \(x = \sum_{b \in B} \alpha_b b\) 。若 \(x \in \mathbf{U}^{\mathrm{H}}\) ，则

\(\alpha_{hb} = \alpha_b\) 对所有 \(b \in \mathbf{B}\) 以及所有 \(h \in \mathbf{H}\) 成立，因此 \(x \in \mathbf{U}'\) ，由此可得 \(\mathbf{U}' = \mathbf{U}^{\mathrm{H}}\) 。

命题 5. --- 设 M 为自由 A-模，k 为整数 \(\geqslant0\) ，P 为 \(\mathbf{TS}^{k}(\mathbf{M})\) 的由 \(\gamma_{k}(\mathbf{M})\) 生成的 A-子模。假设 A 为无限整环。则对每个 \(z\in\mathbf{TS}^{k}(\mathbf{M})\) 存在 \(\alpha\in\mathbf{A}-\{0\}\) 使得 \(\alpha z\in P\) 。

设 K 为 A 的分式域。我们将 \(\mathbf{T}\mathbf{S}^{k}(\mathbf{M})\) 与向量 K-空间 \(\mathbf{V} = \mathbf{T}\mathbf{S}^{k}(\mathbf{M}) \otimes_{\mathbf{A}} \mathbf{K}\) 的一个子 A-模等同（命题 4，以及 II，第 314 页）。我们必须证明这个向量 K-空间由 \(\gamma_{k}(\mathbf{M})\) 生成，即 V 上满足 \(f(\gamma_k(M)) = 0\) 的每个 K-线性形式 f 均为零。设 \((e_i)_{i \in I}\) 为 M 的一组基，并按命题 4 定义 \(e_\nu\) 。对任意 \((\alpha_i) \in A^{(I)}\) ，考虑到 (6)，我们有

\[
0 = f \left(\gamma_ {k} \left(\sum_ {i \in I} \alpha_ {i} e _ {i}\right)\right) = \sum_ {\nu \in \mathbf {N} ^ {(I)}, | \nu | = k} \alpha^ {\nu} f (e _ {\nu}).
\]

由 IV 第 18 页推论 2 可知 \(f(e_{\nu}) = 0\) 对所有 \(\nu \in \mathbf{N}^{(I)}\) 成立，从而 \(f = 0\) 。

